% Generated by scripts/assemble.py; corrections belong in source/corrections.json.
\examhead{36}{华东师范大学 2024 年数学分析试卷}{39}
% p040-b01; source PDF page 40
\blockanchor{p040-b01}
\question{1}(10 分) 计算曲面 $z = e^{x^2 - y^2} + x^2$ 与曲面 $z = e^{x^2 - y^2} - y^2 + 1$ 所围成的封闭区域的体积.\par
% p040-b02; source PDF page 40
\par\Needspace{6\baselineskip}
\blockanchor{p040-b02}
\question{2}(10 分) 设 $f(x)$ 是 $\mathbb{R}$ 上以 $2\pi$ 为周期的连续函数, 以 $a_{n}, b_{n}$ 为其傅里叶级数的系数, 求\par
\nopagebreak[4]
% p040-b03; source PDF page 40
\blockanchor{p040-b03}
\displayformula{F (x) = \frac {1}{\pi} \int_ {- \pi} ^ {\pi} f (t) f (x + t) \mathrm{d} t}
% p040-b04; source PDF page 40
\blockanchor{p040-b04}
的傅里叶系数(用 $a_{n}, b_{n}$ 表示).\par
% p040-b05; source PDF page 40
\blockanchor{p040-b05}
\question{3}(15 分) 设  $\lim_{n\to\infty}\frac{a_{1}+a_{2}+\cdots+a_{n}}{n}=a(-\infty<a<+\infty).$\par
% p040-b06; source PDF page 40
\blockanchor{p040-b06}
(1) 判断 $\lim_{n\to \infty}a_n$ 是否存在, 若存在, 请给出严格证明; 若不存在, 请举例并详细说明.\par
% p040-b07; source PDF page 40
\blockanchor{p040-b07}
(2) 若 $\lim_{n\to \infty}n(a_n - a_{n - 1}) = 0$ 或 $\{a_n\}$ 单调递增, 判断 $\lim_{n\to \infty}a_n$ 是否存在, 若存在, 请给出严格证明; 若不存在, 请举例并给出详细证明.\par
% p040-b08; source PDF page 40
\par\Needspace{6\baselineskip}
\blockanchor{p040-b08}
\question{4}(15 分) 设 $f(x)$ 在 $[0,1]$ 上有连续的导数, $f(0) = 0$ 或 $\int_0^1 f(x)\mathrm{d}x = 0$. 判断不等式\par
\nopagebreak[4]
% p040-b09; source PDF page 40
\blockanchor{p040-b09}
\displayformula{\int_ {0} ^ {1} [ f (x) ] ^ {2} \mathrm{d} x \leq \int_ {0} ^ {1} [ f ^ {\prime} (x) ] ^ {2} \mathrm{d} x}
% p040-b10; source PDF page 40
\blockanchor{p040-b10}
是否成立, 若成立, 请给出严格证明; 若不成立, 请举例并给出详细证明.\par
% p040-b11; source PDF page 40
\blockanchor{p040-b11}
\question{5}(15 分) 讨论 $\sum_{n=1}^{\infty} \frac{(-1)^{n-1}}{[n+(-1)^{n-1}]^p} (p > 0)$ 的敛散性(收敛时需指明是绝对收敛还是条件收敛).\par
% p040-b12; source PDF page 40
\blockanchor{p040-b12}
\question{6}(20 分) 解答如下问题:\par
% p040-b13; source PDF page 40
\blockanchor{p040-b13}
(1) 证明: $\{u_n(x)\} = \left\{\left(1 + \frac{x}{n}\right)^n\right\}$ 在 $x \in [0,1]$ 上一致收敛.\par
% p040-b14; source PDF page 40
\blockanchor{p040-b14}
(2) 证明: $\{v_n(x)\} = \left\{\frac{1}{e^{\frac{x}{n}} + (1 + \frac{x}{n})^n}\right\}$ 在 $x \in [0,1]$ 上一致收敛.\par
% p040-b15; source PDF page 40
\blockanchor{p040-b15}
(3) 求 $\lim_{n\to \infty}\int_0^1\frac{\mathrm{d}x}{e^{\frac{x}{n}} + (1 + \frac{x}{n})^n}$.\par
% p040-b16; source PDF page 40
\blockanchor{p040-b16}
\question{7}(25 分) 设 $f$ 是定义在 $(- \infty, + \infty)$ 上的非常数函数, 请判断下列说法是否正确. 若正确, 请给出严格证明; 若不正确, 请举例并给出详细说明.\par
% p040-b17; source PDF page 40
\blockanchor{p040-b17}
(1) 若 $f$ 是处处不连续的周期函数, 则 $f$ 必有最小正周期.\par
% p040-b18; source PDF page 40
\blockanchor{p040-b18}
(2) 若 $f$ 是处处不连续的周期函数, 则 $f$ 必没有最小正周期.\par
% p040-b19; source PDF page 40
\blockanchor{p040-b19}
(3) 若 $f$ 是周期函数但没有最小正周期, 则 $f$ 必有一列趋于 0 的周期.\par
% p040-b20; source PDF page 40
\blockanchor{p040-b20}
(4) 若 $f$ 是连续的周期函数, 则 $f$ 必有最小正周期.\par
% p040-b21; source PDF page 40
\blockanchor{p040-b21}
\question{8}(20 分) 设 $f$ 在 $(a, +\infty)$ 上有连续的导数且 $\lim_{x \to +\infty} f(x) = A (-\infty < A < +\infty)$.\par
% p040-b22; source PDF page 40
\blockanchor{p040-b22}
(1) 请给出例子详细说明 $\lim_{x \to +\infty} f'(x)$ 未必存在.\par
% p040-b23; source PDF page 40
\blockanchor{p040-b23}
(2) 若 $\lim_{x\to +\infty}f'(x)$ 存在, 求它的值.\par
% p040-b24; source PDF page 40
\blockanchor{p040-b24}
(3) 若  $f(x)$  满足 n 阶可导且  $\lim_{x\to+\infty}f(x)$  与  $\lim_{x\to+\infty}f^{(n)}(x)$  存在. 判断对任意的  $k=1,2,\cdots,n-1$ ,  $\lim_{x\to+\infty}f^{(k)}(x)$  是否存在. 若存在, 请求出极限; 若不存在, 请举例并给出详细说明.\par
% p040-b25; source PDF page 40
\par\Needspace{6\baselineskip}
\blockanchor{p040-b25}
\question{9}(20 分) 设 $f: \mathbb{R}^2 \to \mathbb{R}$ 连续可微, $f(0,0) = 0$ 且 $f_y(0,0) \neq 0$. 考虑迭代:\par
\nopagebreak[4]
% p040-b26; source PDF page 40
\blockanchor{p040-b26}
\displayformula{y _ {n + 1} (x) = y _ {n} (x) - \left(f _ {y} (0, 0)\right) ^ {- 1} f (x, y _ {n} (x)), y _ {0} (x) = 0.}
% p040-b27; source PDF page 40
\blockanchor{p040-b27}
定义 $u_{n + 1}(x) = y_{n + 1}(x) - y_n(x), n \geq 0.$\par
% p041-b00; source PDF page 41
\blockanchor{p041-b00}
(1) 证明: 对任何 $\varepsilon > 0$, 存在 $\delta > 0$, 当 $x \in (-\delta, \delta)$ 时, $|y_n(x)| < \varepsilon$ 对任何 $n \geq 0$ 成立, 进而当 $\delta$ 充分小时, 存在 $0 < \rho < 1$ 使得对任意的 $|x| < \delta$, 有 $|u_{n+1}(x)| \leq \rho |u_n(x)|$.\par
% p041-b01; source PDF page 41
\blockanchor{p041-b01}
(2) 证明: $y(x) = \sum_{n=1}^{\infty} u_n(x)$ 时 $f(x, y) = 0$ 在 (0,0) 附近确定的一个连续的隐函数.\par
